\par\Needspace{6\baselineskip}\medskip
\hypertarget{native-088-article069}{}\pdfbookmark[2]{69.}{article-069}
\begin{center}69.\end{center}
Als Beispiel berechnen wir die Länge und Breite aus der Rektaszension $355^0\,43'\,45''{,}30=\alpha$, der Deklination $-8^0\,47'\,25''=\delta$ und der Schiefe der Ekliptik $23^0\,27'\,59''{,}26=\epsilon$. Es gilt also $45^0+\tfrac12\alpha=222^0\,51'\,52''{,}65$, $45^0-\tfrac12(\epsilon+\delta)=37^0\,39'\,42''{,}87$, $45^0-\tfrac12(\epsilon-\delta)=28^0\,52'\,17''{,}87$. Daraus ergibt sich weiter
\begin{center}
\begin{tabular}{@{}ll@{}}
\makebox[21em][l]{$\log\cos(45^0+\tfrac12\alpha)$\hfill .\,.\,.\,.\,.\,.\,.} & $9{,}865\,0820_n$\\[.3em]
\makebox[21em][l]{$\log\sin(45^0-\tfrac12(\epsilon+\delta))$\hfill .\,.\,.} & $9{,}786\,0418$\\[.3em]
\makebox[21em][l]{$\log\cos(45^0-\tfrac12(\epsilon+\delta))$\hfill .\,.\,.} & $9{,}898\,5222$\\[.3em]
\makebox[21em][l]{$\log\sin(45^0+\tfrac12\alpha)$\hfill .\,.\,.\,.\,.\,.\,.} & $9{,}832\,6803_n$\\[.3em]
\makebox[21em][l]{$\log\sin(45^0-\tfrac12(\epsilon-\delta))$\hfill .\,.\,.} & $9{,}683\,8112$\\[.3em]
\makebox[21em][l]{$\log\cos(45^0-\tfrac12(\epsilon-\delta))$\hfill .\,.\,.} & $9{,}942\,3572$\\[.3em]
\hline
\end{tabular}
\end{center}
\begin{center}
\begin{tabular}{@{}ll@{}}
\makebox[24em][l]{$\log\sin(45^0-\tfrac12 b)\sin\tfrac12(E-l)$\hfill .\,.\,.} & $9{,}651\,1238_n$\\[.3em]
\makebox[24em][l]{$\log\sin(45^0-\tfrac12 b)\cos\tfrac12(E-l)$\hfill .\,.\,.} & $9{,}775\,0375_n$\\[.3em]
\end{tabular}
\end{center}
\noindent woraus sich $\tfrac12(E-l)=216^0\,56'\,5''{,}39$ und $\log\sin(45^0-\tfrac12 b)=9{,}872\,3171$ ergeben.
\begin{center}
\begin{tabular}{@{}ll@{}}
\makebox[24em][l]{$\log\cos(45^0-\tfrac12 b)\sin\tfrac12(E+l)$\hfill .\,.\,.} & $9{,}516\,4915_n$\\[.3em]
\makebox[24em][l]{$\log\cos(45^0-\tfrac12 b)\cos\tfrac12(E+l)$\hfill .\,.\,.} & $9{,}763\,6042_n$\\[.3em]
\end{tabular}
\end{center}
\noindent woraus sich $\tfrac12(E+l)=\grave{2}09^0\,30'\,49''{,}94$ und $\log\cos(45^0-\tfrac12 b)=9{,}823\,9669$ ergeben.

\SourceLeaf{089}Somit ergibt sich $E=426^0\,26'\,55''{,}33$, $l=-7^0\,25'\,15''{,}45$, was auf dasselbe hinausläuft wie $E=66^0\,26'\,55''{,}33$, $l=352^0\,34'\,44''{,}55$. Für den Winkel $45^0-\tfrac12 b$ erhält man aus dem Logarithmus des Sinus $48^0\,10'\,58''{,}11$, aus dem Logarithmus des Kosinus $48^0\,10'\,58''{,}17$ und aus dem Tangens, dessen Logarithmus die Differenz der beiden anderen Logarithmen ist, $48^0\,10'\,58''{,}14$. Daraus folgt $b=-6^0\,21'\,56''{,}28$.

Nach dem anderen Verfahren verläuft die Rechnung wie folgt.
\begin{center}
\begin{tabular}{@{}ll@{}}
\makebox[15em][l]{$\log\operatorname{tang}\delta$\hfill .\,.\,.\,.\,.\,.\,.} & $9{,}189\,3062_n$\\[.3em]
\makebox[15em][l]{$\log\sin\alpha$\hfill .\,.\,.\,.\,.\,.\,.\,.} & $8{,}871\,9792_n$\\[.3em]
\hline
\makebox[15em][l]{$\log\operatorname{tang}\zeta$\hfill .\,.\,.\,.\,.\,.\,.} & $0{,}317\,3270$\\[.3em]
$\zeta=$ & $64^0\,17'\,6''{,}83$\\[.3em]
$\zeta-\epsilon=$ & $40\quad49\quad7{,}57$\\[.3em]
\end{tabular}
\end{center}
\begin{center}
\begin{tabular}{@{}ll@{}}
\makebox[15em][l]{$\mathrm{C.}\log\cos\zeta$\hfill .\,.\,.\,.\,.\,.} & $0{,}362\,6190$\\[.3em]
\makebox[15em][l]{$\log\cos(\zeta-\epsilon)$\hfill .\,.\,.\,.} & $9{,}878\,9703$\\[.3em]
\makebox[15em][l]{$\log\operatorname{tang}\alpha$\hfill .\,.\,.\,.\,.\,.\,.} & $8{,}873\,1869_n$\\[.3em]
\hline
\makebox[15em][l]{$\log\operatorname{tang}l$\hfill .\,.\,.\,.\,.\,.\,.} & $9{,}114\,7762_n$\\[.3em]
$l=$ & $352^0\,34'\,44''{,}50$\\[.3em]
\end{tabular}
\end{center}
\begin{center}
\begin{tabular}{@{}ll@{}}
\makebox[15em][l]{$\log\sin l$\hfill .\,.\,.\,.\,.\,.\,.\,.} & $9{,}111\,1232_n$\\[.3em]
\makebox[15em][l]{$\log\operatorname{tang}(\zeta-\epsilon)$\hfill .\,.\,.} & $9{,}936\,3874$\\[.3em]
\hline
\makebox[15em][l]{$\log\operatorname{tang}b$\hfill .\,.\,.\,.\,.\,.\,.} & $9{,}047\,5106_n$\\[.3em]
$b=$ & $-6^0\,21'\,56''{,}26$\\[.3em]
\end{tabular}
\end{center}
\noindent Zur Bestimmung des Winkels $E$ haben wir zwei Rechenwege.
\begin{center}
\begin{tabular}{@{}ll@{}}
\makebox[15em][l]{$\log\sin\epsilon$\hfill .\,.\,.\,.\,.\,.\,.\,.} & $9{,}600\,1144$\\[.3em]
\makebox[15em][l]{$\log\cos\alpha$\hfill .\,.\,.\,.\,.\,.\,.\,.} & $9{,}998\,7924$\\[.3em]
\makebox[15em][l]{$\mathrm{C.}\log\cos b$\hfill .\,.\,.\,.\,.\,.} & $0{,}002\,6859$\\[.3em]
\hline
\makebox[15em][l]{$\log\cos E$\hfill .\,.\,.\,.\,.\,.\,.} & $9{,}601\,5927$\\[.3em]
\end{tabular}
\end{center}
\begin{center}
\begin{tabular}{@{}ll@{}}
\makebox[15em][l]{$\log\sin\epsilon$\hfill .\,.\,.\,.\,.\,.\,.\,.} & $9{,}600\,1144$\\[.3em]
\makebox[15em][l]{$\log\cos l$\hfill .\,.\,.\,.\,.\,.\,.\,.} & $9{,}996\,3470$\\[.3em]
\makebox[15em][l]{$\mathrm{C.}\log\cos\delta$\hfill .\,.\,.\,.\,.\,.} & $0{,}005\,1313$\\[.3em]
\hline
\makebox[15em][l]{$\log\cos E$\hfill .\,.\,.\,.\,.\,.\,.} & $9{,}601\,5927$\\[.3em]
\end{tabular}
\end{center}
\noindent woraus sich $E=66^0\,26'\,55''{,}35$ ergibt.
